% Some rights are reserved as specified by
% Creative Commons Attribution - NonCommercial version 4.0 licence
% which can be accessed on https://creativecommons.org/licenses/by-nc/4.0/

Copyright 2025 No Board Games
% closing the group opened for avoiding the page break after the ToC

\endgroup

Some rights are reserved as specified by Creative Commons Attribution - NonCommercial version 4.0 licence, which can be accessed on

\url{https://creativecommons.org/licenses/by-nc/4.0/}

\vfill

\begin{minipageparindent}{0.8\linewidth}
    \textbf{No Board Games} is an international group that creates radical board games.
    Check for gameplay videos on our website: \url{https://noboardgames.com}
\end{minipageparindent}%
%\hfill
\begin{minipage}{0.2\linewidth}
    
\includegraphics[width=1\textwidth]{QR_code_NoBG.png}
\end{minipage}

\vfill

\begin{wrapfigure}[6]{R}{0.23\textwidth}
    \vspace{-0.5cm}
    \centering
    
\includegraphics[height=6\baselineskip]{QR_code_BGG_forum.png}
\end{wrapfigure}

\textit{You are reading the version 1.0 of the rulebook, printed in English and Italian. For new
revisions and more translations, check our website. If you want to contribute translations,
contact us at: info@noboardgames.com
The development and discussion on the rules will continue on BoardGameGeek forum
"Rules". Log your games on the forum "Sessions".} \url{https://boardgamegeek.com/boardgame/357812/fully-automated-board-game/forums/}

Printed in Poland by Fabryka Kart, 2025.

\section{ACKNOWLEDGEMENTS}

    \textbf{Design}: Matteo Modena, Ilario Gelmetti
    
    \textbf{Art}: Bernardo "RUPE" Anichini - IG: \texttt{@rupe\_b}
    
    \textbf{Testing}: Ilario Gelmetti, Matteo Modena, Matteo Paolini, Giulio Mazzocchi,Bernardo Anichini,
    Emanuele Croce, Dan Howitt, Luis Gutiérrez, Nicola Fiorillo, Loz, Riccardo Bernori,
    Luna Golino, Michele Bertoldi
    
    \textbf{Communication}: Matteo Paolini, Matteo Modena
    
    \textbf{Video}: Davide Morittu
    
    We want to thank all the backers that trusted us during the crowdfunding campaign, and
    specifically "The Boss" backers: \textbf{Rock Castle}, \textbf{Matteo R.} and \textbf{Comrade Alex}!
    
    Also, this game would not have been possible without the countless online testing
    sessions that were joined by many. Special thanks go to the heroic \textbf{Dan Howitt} and the
    unstoppable \textbf{Luis Gutiérrez}.


\section{BACKGROUND}

    \epigraph{The working class and the employing class have nothing in common.}{\textit{Industrial Workers of the World}}
    
    There will never be peace between the oppressed and the oppressors, and thus between
    master and subordinate. From the first utopian socialists to the present day, a new vision
    of social change has emerged, one based on class struggle. The employers have acquired
    their wealth on the systematic exploitation of the working class, leveraged technology to
    fuel this mechanism, and for avoiding strikes and sabotage. Faced with new technologies,
    capable of speeding up and improving the efficiency of production, some have responded with
    sabotage and destruction, such as the Luddite movement*, and those who
    have claimed all the benefits, such as the Left-wing Accelerationist movement. This game
    encourages you to confront the question of the automation of work, the ever decreasing
    need of wageworkers, and, more generally, with the concept of class struggle. Because
    there will always be class struggle as long as power and resources are not distributed
    equally among people.
    
    \textsl{*The Luddite movement was a class struggle in England during the 19th century, whereby
    workers struggled against their labours being replaced with machinery.}


\section{OVERVIEW}

    Take the role of a greedy boss (representing the Bourgeoisie) or suffer the burden of a
    worker (representing the Proletariat).
    
    As a \textbf{worker}, you must negotiate and accept Contracts from the bosses, to collect enough
    coins to pay your Rent and \textbf{survive}. But only a collective and coordinated \textbf{strike} will bring
    victory to the workers.
    
    As a \textbf{boss}, you control the production, the money, and therefore the workers. You must
    offer the favoured Contracts and \textbf{accumulate the greatest profit} before your competitor
    does.
    
    The Fully Automated Board Game is a fusion of negotiation, hidden roles, bluffing, voting,
    and class struggle.


\section{COMPONENTS}


    \subsubsection{PHASE MARKER TOKEN}
        \begin{wrapfigure}[7]{R}{0.23\textwidth}
            \vspace{-0.5cm}
            \centering
            \includegraphics[height=6\baselineskip]{token_phase.jpg}
        \end{wrapfigure}
        In this box, we decided to not include any component representing
        the \faRubleSign\ money. Instead, we propose
        you to use real coins. This will make
        the gameplay feel more "real" and
        allows us to reduce the materials,
        thus making the box smaller, more
        portable, and its fabrication less
        polluting. For the same reason, this
        game has been fabricated with the
        eco-friendly materials and methods of Fabryka Kart in Poland.
        At the end of each game session,
        we encourage you to take all the
        real coins used as in-game currency and make a donation to a
        union or a social project, in the real
        world! Consider this a way of helping your local social community.
        Reach out, maybe we can recommend some entity!

    \subsubsection{WHAT'S IN THE BOX}
        \begin{multicols}{2}
            \begin{itemize}
                \item 2 Boss cards - Bourgeoisie
            (blue card with a blue pig symbol)
            
                \item 10 Worker character cards:
            
                \begin{itemize}
                    \item 4 Proletarians - Virtual pauper
            (red card with a red flag symbol)
            
                    \item 3 Startupper - Petite bourgeoisie
            (red card with a blue pig symbol)
            
                    \item 2 Scab - Workaholic (red card with a
            symbol of a mug "I \textless3 Mondays")
            
                    \item 1 Luddite - Saboteur
            (red card with a bomb symbol)
            \end{itemize}
            
                \item 10 Robot cards (blue and red sides)
                \item 17 Event cards
                \item 6 Contracts cards
                \item 1 Union Fund card
                \item 1 Public Work card
                \item 1 Mayor card
                \item 1 Security Guard card
                \item 6 Phase cards
                \item 8 Personal tokens for workers
                \item 1 Phase marker token
            \end{itemize}
        \end{multicols}
